\section{Introduction}
\label{sec:intro}

Forty-nine of the fifty U.S.\ states have bicameral legislatures, with a lower
house and an upper senate chamber. Many state constitutions require or encourage
senate districts to be composed of whole house districts---a \emph{nesting}
requirement that ensures geographic coherence between the two chambers.
California's constitution, for example, requires that ``each Senate district
shall consist of two complete, contiguous Assembly districts.''
Iowa's statute requires that senate districts shall consist of whole house
districts wherever practicable.

Nesting imposes a geometric constraint: if senate district $S_j$ must be
composed of whole house districts, the house map must first be drawn, and the
senate map derived from it. This creates an algorithmic dependency: the house
map must be drawn to enable a valid senate map, not merely to optimise its
own compactness. A house map that is individually optimal may fragment into
senate districts that are non-compact or non-contiguous.

Companion paper B.13 \citep{dellaLibera2026nestSection} introduced the
\emph{NestSection} algorithm: given house count $H$ and senate count $S$,
construct a geographic \emph{spine} at the level of $g = \gcd(H, S)$
sub-regions, then subdivide each spine region into $H/g$ house districts
and $S/g$ senate districts. NestSection guarantees that senate districts are
exact unions of house districts while jointly optimising both chambers' compactness.

This paper applies NestSection to all 49 bicameral state legislatures.
Our contributions are:
\begin{enumerate}
  \item A complete catalogue of $H:S$ ratios for all 49 bicameral states,
    classified by NestSection compatibility ($\gcd(H,S) > 1$ vs.\ $= 1$).
  \item Application of NestSection to all 42 compatible states, producing
    joint house-senate maps with verified nesting.
  \item Quantification of the compactness penalty from nesting: mean $+2$--$3\%$
    in edge-cut relative to independently optimised maps.
  \item Measurement of the partisan effect of nesting constraints: $\leq 0.3$
    expected seat shift across compatible states.
  \item A classification of the 7 incompatible states ($\gcd = 1$) and
    a discussion of constitutional alternatives.
\end{enumerate}

\paragraph{Organization.}
Section~\ref{sec:nestsection} recaps the NestSection algorithm.
Section~\ref{sec:analysis} performs the 50-state ratio analysis.
Section~\ref{sec:results} reports NestSection results for the 42 compatible states.
Section~\ref{sec:penalty} quantifies the compactness penalty.
Section~\ref{sec:partisan} measures the partisan effect.
Section~\ref{sec:conclusion} concludes.
